% Continuously revised edition; last source revision 24 September 2026.
% Every discrepancy from the printed article is explained in jones1980_editorial_notes.md;
% a dated log of revisions and checks is ../EDITORIAL_NOTES.md.
% Corrected transcription of James P. Jones, Bull. Amer. Math. Soc. (N.S.)
% 3 (1980), 859--862. DOI: 10.1090/S0273-0979-1980-14832-6.
% Prepared from the supplied scan and Mathpix OCR. See jones1980_editorial_notes.md.
% Compile twice with: pdflatex jones1980_corrected.tex
\documentclass[11pt,letterpaper]{article}
\usepackage[T1]{fontenc}
\usepackage[utf8]{inputenc}
\usepackage{lmodern}
\usepackage[margin=0.82in,headheight=14pt,marginparwidth=0.55in,marginparsep=4pt]{geometry}
\usepackage{amsmath,amssymb,amsthm}
\usepackage{microtype}
\usepackage{booktabs}
\usepackage{fancyhdr}
\usepackage{marginnote}
\usepackage[hidelinks]{hyperref}
% Original pagination: \origpage{n} puts the number of the Bulletin page
% that begins on the current line into the margin (line accuracy).
\newcommand{\origpage}[1]{\marginnote{\footnotesize\textbf{[#1]}}}
\let\origpagemark\origpage
\hypersetup{pdftitle={Undecidable Diophantine Equations -- corrected transcription},
 pdfauthor={James P. Jones},pdfsubject={Corrected transcription of the 1980 article}}
\newtheorem{theorem}{Theorem}
\newcommand{\Matijasevic}{Matijasevi\v{c}}
\pagestyle{fancy}
\fancyhf{}
\fancyhead[L]{\small JAMES P. JONES}
\fancyhead[R]{\small UNDECIDABLE DIOPHANTINE EQUATIONS}
\fancyfoot[C]{\thepage}
\renewcommand{\headrulewidth}{0pt}
\setlength{\parindent}{1.3em}
\setlength{\parskip}{2pt}
\setlength{\emergencystretch}{1.5em}
\setlength{\jot}{5pt}
\title{\vspace{-1.5em}\Large\bfseries UNDECIDABLE DIOPHANTINE EQUATIONS}
\author{\normalsize JAMES P. JONES}
\date{}
\begin{document}
\maketitle
\thispagestyle{plain}
\vspace{-1.8em}
\begin{center}
\small \textit{Bulletin (New Series) of the American Mathematical Society}\\
Volume 3, Number 2, September 1980, pp.\ 859--862\\
\href{https://doi.org/10.1090/S0273-0979-1980-14832-6}{doi:10.1090/S0273-0979-1980-14832-6}
\end{center}
\vspace{-0.5em}
\begin{quote}\footnotesize
Corrected transcription, continuously revised: every discrepancy between the printed announcement and this edition is explained in \texttt{Papers/1980/}\allowbreak\texttt{jones1980\_editorial\_notes.md}, and a dated log of revisions and checks is \texttt{Papers/EDITORIAL\_NOTES.md}. Mathematical and bibliographic
emendations are documented in the accompanying editorial notes. The results
and bounds are those announced in the 1980 article, not a survey of later work.
The margin numbers [859]--[862] mark the lines on which the pages of the
original printing begin.
\end{quote}

% Original p. 859 (title page). OCR: distinguish Greek nu (arity) from Latin v (index).
\origpagemark{859}In 1900 Hilbert asked for an algorithm to decide the solvability of all
diophantine equations, $P(x_1,\ldots,x_\nu)=0$, where $P$ is a polynomial
with integer coefficients. In special cases of Hilbert's tenth problem,
such algorithms are known. Siegel~\cite{Siegel} gives an algorithm for all
polynomials $P(x_1,\ldots,x_\nu)$ of degree $\leqslant 2$. From the work of
A.~Baker~\cite{Baker} we know that there is also a decision procedure for
the case of homogeneous polynomials in two variables, $P(x,y)=c$.

An essential step toward the eventual negative solution of the entire
(unrestricted) form of Hilbert's tenth problem was taken in 1961 by Julia
Robinson, Martin Davis and Hilary Putnam~\cite{DPR}. They proved that every
recursively enumerable set $W$ can be represented in exponential
diophantine form
\[
 x\in W\ \Longleftrightarrow\
 \exists x_1,\ldots,x_\mu\;
 P(x,x_1,\ldots,x_\mu,2^{x_1},\ldots,2^{x_\mu})=0,
\]
where $P$ is a polynomial with integer coefficients and $x_1,\ldots,x_\mu$
range over positive integers.

In 1970 Ju.~V.~\Matijasevic~\cite{Mat70} proved that the exponential
relation $y=2^x$ is diophantine and hence that every r.e.\ set $W$ can be
represented in polynomial (diophantine) form
\[
 x\in W\ \Longleftrightarrow\
 \exists x_1,\ldots,x_\nu\;
 P(x,x_1,\ldots,x_\nu)=0,
\]
where the unknowns $x_1,\ldots,x_\nu$ range over positive integers.

Since there exist r.e.\ nonrecursive sets, \Matijasevic's theorem implies
the undecidability of Hilbert's tenth problem. There is no algorithm to
decide whether an arbitrary diophantine equation has a solution.

\Matijasevic's theorem implies also the existence of particular undecidable
families of diophantine equations. In fact there must exist universal
diophantine equations, polynomial analogues of the universal Turing
machine. This follows from the well-known fact that the r.e.\ sets
$W_1,W_2,\ldots$ can be listed in such a way that the binary relation
$x\in W_v$ is r.e.\ Hence by~\cite{DPR,Mat70} there exists a universal
polynomial $U(x,v,x_1,\ldots,x_\nu)$ with the property
\[
 x\in W_v\ \Longleftrightarrow\
 \exists x_1,\ldots,x_\nu\;
 U(x,v,x_1,\ldots,x_\nu)=0.
\]
Thus a single polynomial, of a fixed degree and a fixed number of
unknowns, can define every r.e.\ set by mere change of a parameter $v$.
The existence of such\origpage{860} a universal polynomial follows immediately
from~\cite{DPR,Mat70}. However, the construction of an actual example is
a different matter. The first example to be specifically written down
is given in~\cite{Jones78}. More recently the author has constructed
further examples. These are given below, without proof. (The proof uses
the code idea of~\cite{MR75} together with new methods of \Matijasevic\
based on E.~Kummer's theorem about divisibility of binomial coefficients
by prime powers.)

\newpage
% Original p. 860 begins at the \origpage{860} mark above, inside the last
% paragraph; the forced break here keeps Theorem 1 on one page. OCR: q=b^{560} -> q=b^{5^{60}}.
\noindent\textit{Editorial index convention.} In Theorems 1--3 the triples
$\langle z,u,y\rangle$ are admissible polynomial codes: the triples constructed
in (4.1) of Jones's 1982 \textit{Universal Diophantine Equation}, with the
normalization specified in its \S3. This is the class of indices for which
the common equivalences are proved. The precise coefficient convention is
recorded in the editorial notes (entry 1980-08); an arbitrary unrelated
positive triple is not automatically such a code.

\begin{theorem}
In order that $x\in W_{\langle z,u,y\rangle}$, it is necessary and
sufficient that the following system of equations have a solution in
positive integers.
\end{theorem}
\[
\begin{gathered}
 e\ell g^2+\alpha=(b-xy)q^2,\qquad q=b^{5^{60}},\\
 \lambda+q^4=1+\lambda b^5,\qquad \theta+2z=b^5,\\
 \ell=u+t\theta,\qquad e=y+m\theta,\qquad b=2^w,\\[2pt]
 \binom{g+q^3-1-b\ell+\ell}{g}
 \binom{\theta\lambda+e+\ell q^2}{\theta\lambda}\\[2pt]
 {}\cdot
 \binom{b^5q-2q+2(e-z\lambda)(1+xb^5+g)^4+b^5\lambda(1+q^4)}{b^5q-2q}
 =2\eta+1.
\end{gathered}
\]
Here we consider the r.e.\ sets to be indexed by three indices, $z,u,y$,
instead of the usual one. This is an inessential restriction. To use
indexing with a single parameter $v$, one need only add to our equations
the equation
\[
 v=\bigl((zuy)^2+u\bigr)^2+y
\]
and regard $z,u,y$ as unknowns instead of parameters.\footnote{Editorial
clarification: this gives a suitable effective one-parameter indexing;
it need not coincide numerically with a previously chosen indexing.
The displayed coding is injective on positive-integer triples.}

As stated above, our system of equations has 12 unknowns, specifically
$b,e,g,\ell,m,q,t,w,\alpha,\eta,\theta,\lambda$, and four parameters,
$z,u,y,x$. An exponential equation, $b=2^w$, appears, and the last line
is a product of three binomial coefficients. In the next system the
number of binomial coefficients is reduced to one.

\begin{theorem}
In order that $x\in W_{\langle z,u,y\rangle}$, it is necessary and
sufficient that the following system of equations have a solution in
positive integers.
\end{theorem}
\[
\begin{gathered}
 e\ell g^2+\alpha=(b-xy)q^2,\qquad q=b^{5^{60}},\\
 \lambda+q^4=1+\lambda b^5,\qquad \theta+2z=b^5,\\
 \ell=u+t\theta,\qquad e=y+m\theta,\qquad b=2^w,\qquad n=q^{16},\\
\begin{aligned}
 r={}&\Bigl[g+eq^3+\ell q^5\\
     &\quad+\bigl(2(e-z\lambda)(1+xb^5+g)^4
                   +\lambda b^5+\lambda b^5q^4\bigr)q^7\Bigr](n^2-n)\\
     &+\Bigl[q^3-b\ell+\ell+\theta\lambda q^3+(b^5-2)q^8\Bigr](n^2-1),
\end{aligned}\\[2pt]
 % OCR: n n^2 -> eta n^2. It is NOT eta 2^n.
 \eta n^2=\binom{2r}{r}.
\end{gathered}
\]
Here there are fourteen unknowns,
$b,e,g,\ell,m,n,q,r,t,w,\alpha,\eta,\theta,\lambda$, and the four parameters
$x,z,u,y$. Only one binomial coefficient and one exponential\origpage{861} function
appear. Next these are eliminated so that we obtain a system of purely
polynomial equations.

\newpage
% Original p. 861 begins at the \origpage{861} mark above; the forced break
% here keeps Theorem 3 on one page.
\enlargethispage{4\baselineskip}
\begin{theorem}
In order that $x\in W_{\langle z,u,y\rangle}$, it is necessary and
sufficient that the following system of equations have a solution in
positive integers.\footnote{Editorial emendation: the equation for $d$
contains $4a\gamma$, with Latin $a$, not $4\alpha\gamma$ as printed in the
1980 article. The latter is an original typographical error, not an OCR
error; see the accompanying notes for corroboration and the Pell-equation
congruence that supports the correction.}
\end{theorem}
\[
\begin{gathered}
 e\ell g^2+\alpha=(b-xy)q^2,\qquad q=b^{5^{60}},\\
 \lambda+q^4=1+\lambda b^5,\qquad \theta+2z=b^5,\\
 \ell=u+t\theta,\qquad e=y+m\theta,\qquad n=q^{16},\\
\begin{aligned}
 r={}&\Bigl[g+eq^3+\ell q^5\\
     &\quad+\bigl(2(e-z\lambda)(1+xb^5+g)^4
                   +\lambda b^5+\lambda b^5q^4\bigr)q^7\Bigr](n^2-n)\\
     &+\Bigl[q^3-b\ell+\ell+\theta\lambda q^3+(b^5-2)q^8\Bigr](n^2-1),
\end{aligned}\\
 p=2ws^2r^2n^6,\qquad p^2k^2-k^2+1=\tau^2,\\
 4(c-ksn^2)^2+\eta=k^2,\\
 k=r+1+hp-h,\qquad a=(wn^2+1)rsn^2,\\
 c=2r+1+\varphi,\qquad d=bw+ca-2c+4a\gamma-5\gamma,\\
 d^2=(a^2-1)c^2+1,\qquad f^2=(a^2-1)i^2c^4+1,\\
 (d+of)^2=\Bigl(\bigl(a+f^2(d^2-a)\bigr)^2-1\Bigr)(2r+1+jc)^2+1.
\end{gathered}
\]
The equations of Theorem~3 have twenty-eight unknowns,
\[
 a,b,c,d,e,f,g,h,i,j,k,\ell,m,n,o,p,q,r,s,t,w,
 \alpha,\gamma,\eta,\theta,\lambda,\tau,\varphi.
\]
The degree is $5^{60}$; however, the equation $q=b^{5^{60}}$ can be
replaced by certain others of low degree. In fact, by introducing some
30 additional unknowns and new equations one can reduce the degree of
the system to 2. Then, by transposing terms to one side and summing
squares, one can construct a universal diophantine equation in
58 unknowns and of degree 4.

Alternatively, one may try instead to reduce the total number of
unknowns, $\nu$. In~\cite{MR75} Julia Robinson and Ju.~\Matijasevic\ showed
that $\nu$ can be reduced universally to 13. More recently
\Matijasevic~\cite{Mat77} has improved this to $\nu=9$. The corresponding
value of the degree, $\delta$, is, however, very large. The following
table gives various simultaneous possibilities for $\delta$ and $\nu$,
sufficient for a universal equation.

\begin{theorem}
The following universal constructions are reported. The large degrees
written in scientific notation are approximate sizes, not exact
integer degree bounds.
\end{theorem}
\begin{center}
\renewcommand{\arraystretch}{1.13}
\begin{tabular}{@{}rr@{\qquad\qquad}rr@{}}
\toprule
$\nu$ & $\delta$ & $\nu$ & $\delta$\\
\midrule
58 & 4  & 21 & 96\\
38 & 8  & 19 & 2668\\
32 & 12 & 14 & $\approx 2.0\times10^5$\\
29 & 16 & 13 & $\approx 6.6\times10^{43}$\\
28 & 20 & 12 & $\approx 1.3\times10^{44}$\\
26 & 24 & 11 & $\approx 4.6\times10^{44}$\\
25 & 28 & 10 & $\approx 8.6\times10^{44}$\\
24 & 36 & 9  & $D_9$\\
\bottomrule
\end{tabular}
\end{center}
\noindent Here the nine-unknown degree is the expression reported in the
later full derivation (Jones, 1982, \S3):
\[
 D_9=47216\cdot5^{58}+9728
     \approx1.638134073\times10^{45}<1.7\times10^{45}.
\]
\noindent The printed $1.6\times10^{45}$ rounds this expression downward.
The editorial notes (entry 1980-18) distinguish evaluation of this exact expression
from a fresh derivation of the complete relation-combining degree.


% Original p. 862 begins here. New pagination continues without a forced page break.
\origpage{862}A different measure of size and complexity of a system of diophantine
equations is the total number of indicated arithmetical operations, $o$,
i.e.\ the number of additions, subtractions and multiplications necessary
to evaluate or determine the correctness of a proposed solution. For the
equations of Theorem~3, modified to do away with $q=b^{5^{60}}$, it can be
shown that $o=100$. This number, $o$, can be given an interesting
interpretation as a proof-theoretic complexity bound.

Via G\"odel numbering, the theorems of an effectively axiomatizable theory
$T$ become, in effect, an r.e.\ set, and the search for proofs becomes the
search for solutions of a diophantine equation. A formal deduction is effectively recoverable from such a solution,
but this recoverability does not assert preservation of proof length
or computational complexity. Hence we have

\begin{theorem}
For any effectively axiomatizable theory $T$ and any proposition $P$, if
$P$ has a proof in $T$, then $P$ has another proof consisting of
100 additions and multiplications of integers.\footnote{Editorial
clarification: ``another proof'' here means an arithmetic certificate
under a fixed effective encoding of provability, not necessarily a
100-step derivation in the original proof calculus of $T$. The count is
of the indicated arithmetic operations on whole integers in the modified
system; it is not a bound on the lengths of the integers, the number of
bit operations, or the work of finding a certificate. The preceding
operation-count definition includes subtraction. The count 100 is
reproduced by counting every indicated sign of Theorem~3 once (powers not
counted); that reading, and certificates under a calculator model, are
examined in \texttt{jones1980\_theorem5\_operations.pdf} in this directory.}
\end{theorem}

For the universal equations in~\cite{Jones78} the value of $o$ was 243.
The bound obtained here is $o=100$. As with $\nu$ and $\delta$, it would be
interesting to know the minimum value of $o$.

\begin{thebibliography}{7}
\small
\bibitem{Baker}
A.~Baker, \textit{Contributions to the theory of diophantine equations. I},
Philos.\ Trans.\ Roy.\ Soc.\ London Ser.~A \textbf{263} (1968), 173--191.

\bibitem{DPR}
Martin Davis, Hilary Putnam and Julia Robinson,
\textit{The decision problem for exponential diophantine equations},
Ann.\ of Math.\ (2) \textbf{74} (1961), 425--436.

\bibitem{Jones78}
J.~P.~Jones, \textit{Three universal representations of recursively
enumerable sets}, J.\ Symbolic Logic \textbf{43} (1978), 335--351.

\bibitem{Mat70}
Ju.~V.~\Matijasevic, \textit{Enumerable sets are diophantine},
Dokl.\ Akad.\ Nauk SSSR \textbf{191} (1970), 279--282.
English transl.: Soviet Math.\ Doklady \textbf{11} (1970), 354--357.

\bibitem{Mat77}
Ju.~V.~\Matijasevic, \textit{Some purely mathematical results inspired by
mathematical logic}, in \textit{Logic, Foundations of Mathematics, and
Computability Theory} (Robert E.~Butts and Jaakko Hintikka, eds.),
Reidel, Dordrecht, Holland, 1977, pp.\ 121--127.

\bibitem{MR75}
Ju.~V.~\Matijasevic\ and Julia Robinson,
\textit{Reduction of an arbitrary diophantine equation to one in
13 unknowns}, Acta Arith.\ \textbf{27} (1975), 521--553.

\bibitem{Siegel}
Carl L.~Siegel, \textit{Zur Theorie der quadratischen Formen},
Nachr.\ Akad.\ Wiss.\ G\"ottingen Math.-Phys.\ Kl.~II 1972, 21--46.
\end{thebibliography}

\vspace{0.5em}
\noindent\small\textsc{Department of Mathematics and Statistics,
University of Calgary, Calgary, Alberta, Canada, T2N 1N4}

\vspace{1em}
\noindent\footnotesize
Received by the editors March 5, 1980.\\
\textit{AMS (MOS) subject classifications (1970).}
Primary 02F25, 10N05; Secondary 10B15.\\
\copyright\ 1980 American Mathematical Society. Original publication code:
\texttt{0002-9904/80/0000-0406/\$02.00}.
\end{document}
